% Lecture 10: Backends for mobile: serverless, VPS, and Go
% Compile: latexmk -xelatex lecture10.tex
\documentclass[10pt,aspectratio=169,t]{beamer}
\usetheme{upbminimal}

\course{Mobile and Embedded Computing}
\decklabel{Lecture 10}
\title{Backends for mobile}
\subtitle{Serverless, a VPS, and Go}
\author{Dinu-Ștefan Rusu}
\institute{FILS, Universitatea Politehnica București}
\date{Spring 2027}

\begin{document}

\titleframe

\objectivesframe{
  \cell[\faIcon{code-branch}]{Split client and server}{What must run on a machine you control, and why the client is untrusted input.}
  \cell[\faCloud]{Choose serverless or a VPS}{Cold starts, scaling and cost shape, and the case where each one wins.}
  \cell[\faIcon{plug}]{Design an API for mobile}{Pagination, versioning, conditional requests, compression and rate limits.}
  \cell[\faServer]{Ship a Go backend}{A minimal Docker image, a systemd unit, environment variables, health checks and logs.}
}

\divider{Part 1 \textperiodcentered\ Server and client}{Where does the code run?}[The same feature costs different things on either side of the wire]

\begin{frame}[eyebrow=Server and client]{Server-side and client-side execution}
  \begin{columns}[T]
    \begin{column}{0.6\textwidth}
      \begin{itemize}
        \item Client-side code runs on the phone: rendering, local state, caching, and the optimistic update that makes a form feel instant.
        \item Server-side code runs on a machine you control: secrets, durable storage, authorization, billing, and anything that must be identical for every user.
        \item One request and one response cross the network between them. On mobile that round trip costs roughly 40–200 ms on a good connection, seconds on a bad one, and nothing at all inside a tunnel.
      \end{itemize}
    \end{column}
    \begin{column}{0.36\textwidth}
      \kv{Request lifecycle}{}
      \begin{enumerate}
        \item Client builds the request and holds the token.
        \item Network: a handshake, then the round trip.
        \item Server: authenticate, authorize, execute.
        \item Response: a status code and a body.
      \end{enumerate}
    \end{column}
  \end{columns}
  \begin{takeaway}{Put work on the client for latency.}
    Put it on the server when the answer has to be correct for every user.
  \end{takeaway}
  \note{Ask the room where they would put: a password strength check (both), a price calculation (server), a dark-mode toggle (client), whether a coupon is still valid (server, always). The client copy is a courtesy; the server copy is the rule.}
\end{frame}

\begin{frame}[eyebrow=Server and client]{The client is untrusted input}
  \begin{columns}[T]
    \begin{column}{0.58\textwidth}
      \begin{itemize}
        \item Everything that arrives from a client is user input, including fields your own app filled in.
        \item An Android package (APK) can be unpacked and release Dart code inspected. There is no secret you can ship inside a client.
        \item Client-side validation is a courtesy: an instant warning, a disabled button, no wasted round trip.
        \item The server repeats every check it cares about; it is the only copy nobody else can edit.
      \end{itemize}
    \end{column}
    \begin{column}{0.38\textwidth}
      \begin{hint}[Server-side only]
        API keys, and anything signed with one. Who may read or write which record. Prices, totals, discounts, balances. Rate limits and quotas.
      \end{hint}
    \end{column}
  \end{columns}
  \begin{takeaway}{Validate on the client for speed, on the server for correctness.}
    Any rule enforced only in the app can be bypassed.
  \end{takeaway}
  \note{An Application Programming Interface (API) key is the clearest example: ship it inside the app and every user has it within a day of release. The rule is simple: if leaking it would hurt, it does not belong on the client.}
\end{frame}

\divider{Part 2 \textperiodcentered\ Where the server runs}{Serverless vs a VPS}[Two cost shapes, two failure modes, and two different sets of operational work]

\begin{frame}[eyebrow=Serverless]{Serverless: functions as a service}
  \begin{columns}[T]
    \begin{column}{0.6\textwidth}
      \begin{itemize}
        \item Functions as a Service (FaaS) means you deploy a function, not a machine. The platform provisions, scales, patches, and terminates Transport Layer Security (TLS) connections for you.
        \item It creates instances on demand and destroys them when traffic stops, scaling all the way down to zero.
        \item Billing is per invocation plus gigabyte-seconds of memory, so no traffic means no bill.
        \item AWS Lambda, Google Cloud Run functions, Azure Functions, and Cloudflare Workers all follow this model.
      \end{itemize}
    \end{column}
    \begin{column}{0.36\textwidth}
      \kv{The shape of it}{}
      \begin{enumerate}
        \item Your handler: one function, one job.
        \item The platform: provision, scale, patch, terminate TLS.
        \item 0 to n instances, killed when idle.
      \end{enumerate}
    \end{column}
  \end{columns}
  \begin{takeaway}{You stop operating servers.}
    In exchange you work within the platform's limits.
  \end{takeaway}
  \note{Ask who has already used a FaaS platform without knowing the term. Most students have, through Cloud Functions for Firebase or a hosting provider's edge functions. The unfamiliar part is usually the billing model, not the code.}
\end{frame}

\begin{frame}[eyebrow=Serverless]{Cold starts: the latency of the first request}
  \vskip 1mm
  {\usebeamerfont{small}\begin{tabular}{@{}lll@{}}
    \toprule
    \thead{Case} & \thead{Startup} & \thead{Why} \\
    \midrule
    Warm instance & 0 ms & the container is already running \\
    Go or Rust binary & 10–50 ms & nothing to boot: one static executable \\
    JVM, .NET, large Node graphs & 0.3–2 s & runtime start, plus every import \\
    \bottomrule
  \end{tabular}}
  \vskip 4mm
  \begin{itemize}
    \item A cold start is the platform building a new instance because no warm one was free: pull the image, start the runtime, run top-level code, then finally the handler.
    \item It lands on the first user after a period of no traffic, and again on every traffic spike.
  \end{itemize}
  \muted{Minimum instances, or provisioned concurrency, mean paying for idle capacity, the exact cost serverless was meant to remove. A smaller package and a fast-starting runtime cost nothing extra, which is one reason for the Go section later in this lecture.}
  \note{Numbers are orders of magnitude, not benchmarks. Tell students to measure their own p99 on the platform they picked. Cold-start cost is unusually a language choice: almost nothing else about a small backend depends on the runtime.}
\end{frame}

\begin{frame}[eyebrow=VPS]{VPS: renting the whole machine}
  \begin{columns}[T]
    \begin{column}{0.6\textwidth}
      \begin{itemize}
        \item A Virtual Private Server (VPS) is a virtual machine you rent from a provider such as Hetzner Cloud, DigitalOcean, or OVH: root access, one IP address, and nothing else.
        \item Everything above the hypervisor is yours: operating system updates, firewall rules, TLS certificates, process supervision, backups, monitoring.
        \item Billing is hourly or monthly for the machine, not for the requests it serves. An idle VPS costs what a busy one costs.
      \end{itemize}
    \end{column}
    \begin{column}{0.36\textwidth}
      \kv{The shape of it}{}
      \begin{enumerate}
        \item Your binary: systemd or Docker restarts it.
        \item The OS you patch: firewall, TLS certificates.
        \item The VM you rent: CPU, RAM, disk, a public IP.
      \end{enumerate}
    \end{column}
  \end{columns}
  \begin{takeaway}{Nobody patches it for you.}
    In exchange, no platform limits what you can run.
  \end{takeaway}
  \note{A floating IP, an address you can repoint at a fresh box during an upgrade, is worth mentioning: it allows replacing a machine without changing DNS. Hetzner and DigitalOcean both offer it.}
\end{frame}

\begin{frame}[eyebrow=Serverless vs VPS]{The same backend, two cost shapes}
  {\fontsize{8}{9.5}\selectfont\renewcommand\arraystretch{1.0}
  \begin{tabular}{@{}lp{42mm}p{42mm}@{}}
    \toprule
    \thead{Aspect} & \thead{Serverless (FaaS)} & \thead{VPS} \\
    \midrule
    Scaling & automatic, zero to peak and back & you resize the box yourself \\
    Cost when idle & nothing & the full monthly price \\
    Cost under load & per invocation and memory time & flat, until you outgrow the machine \\
    First request & a cold start, unless paid to avoid it & already warm \\
    You operate & your handler & OS, runtime, TLS, backups, monitoring \\
    Long-running work & capped by the platform timeout & runs as long as you let it \\
    Wins when & spiky, small, or unknown traffic & steady traffic, or state on disk \\
    \bottomrule
  \end{tabular}}
  \vskip 1mm
  \begin{takeaway}{Pick by cost shape and operational load.}
    Spiky or unknown favors serverless. Steady and predictable favors one small VPS.
  \end{takeaway}
  \note{This table is the one slide worth photographing. Every project team will face this choice for the backend in Lab 5 and Lab 6, and the honest answer for a semester project is usually one small VPS.}
\end{frame}

\begin{frame}[eyebrow=Serverless]{Writing code that survives a FaaS runtime}
  \begin{grid}[3]
    \cell[\faIcon{box}]{Ship a small package}{Every megabyte and import adds cold-start time. Trim dependencies first.}
    \cell[\faIcon{sync-alt}]{Initialize outside the handler}{Database pools and HTTP clients belong at module scope, reused by warm invocations.}
    \cell[\faDatabase]{Keep the handler stateless}{Local disk and memory vanish between calls. State belongs in a database or a cache.}
    \cell[\faIcon{redo}]{Make it idempotent}{The platform retries on its own. A duplicate must not charge a card twice.}
    \cell[\faClock]{Respect the timeout}{Work that can exceed it belongs on a queue with a job record, not a longer function.}
    \cell[\faIcon{cogs}]{Know your concurrency}{One provider serves one request per instance; another serves many at once.}
  \end{grid}
  \note{Every one of these is a bug students will eventually hit: a handler that opens a new database connection per call, or stores a counter in a global variable that resets on the next cold start.}
\end{frame}

\begin{frame}[eyebrow=VPS]{Running a box you will not regret}
  \vskip 2mm
  \begin{grid}[3]
    \cell[\faIcon{terminal}]{Rebuild it from a script}{A Dockerfile plus one provisioning script. The machine should be reproducible, never hand-tuned.}
    \cell[\faIcon{sync-alt}]{Run it as a service}{systemd or Docker with a restart policy. A crash at 3 a.m. must recover without you.}
    \cell[\faLock]{Lock down access}{Key-only SSH, no root login, a firewall open only on ports 80 and 443, automatic security updates.}
    \cell[\faIcon{shield-alt}]{Terminate TLS properly}{Caddy or nginx in front, with automatic certificates. Never serve an API over plain HTTP.}
    \cell[\faIcon{save}]{Test the restore}{Snapshots plus an off-box copy of the database. A backup you never restored is unverified.}
    \cell[\faIcon{bell}]{Make it page you}{An uptime check, log shipping, a disk-full alert. A bare VPS does not tell you it went down.}
  \end{grid}
  \note{Ask who has ever restored a backup they took. Almost nobody has, which is the point of the fifth cell: an untested backup is a hope, not a plan.}
\end{frame}

\divider{Part 3 \textperiodcentered\ Go backends}{A backend that fits in one binary}[One static binary, milliseconds to start, deployable anywhere]

\begin{frame}[fragile,eyebrow=Go backends]{Why Go for a small mobile backend}
  \begin{columns}[T]
    \begin{column}{0.6\textwidth}
      \lead{One static binary.}{One executable: no runtime, no dependency folder.}
      \lead{It starts in milliseconds.}{Cold starts on a scale-to-zero platform stay short.}
      \lead{It cross-compiles from your laptop.}{One variable retargets it: a server or a Pi.}
      \lead{The standard library is already the server.}{net/http, encoding/json, goroutines (Lecture 2).}
    \end{column}
    \begin{column}{0.36\textwidth}
\begin{shell}
CGO_ENABLED=0 GOOS=linux \
  GOARCH=arm64 \
  go build -o server
scp server \
  root@vps:/opt/api/
ssh root@vps \
  systemctl restart api
\end{shell}
    \end{column}
  \end{columns}
  \begin{takeaway}{A small binary that starts fast.}
    Those are the properties a mobile client's backend needs.
  \end{takeaway}
  \note{Cross-compilation is the detail students forget: the same laptop that runs Flutter builds a Linux server binary and a Raspberry Pi binary with one environment variable each, no cross-compiler toolchain to install.}
\end{frame}

\begin{frame}[fragile,eyebrow=Go backends]{A handler, and the Dart call that hits it}
  \begin{columns}[T]
    \begin{column}{0.52\textwidth}
      {\bfseries Go: standard library only}\par
\begin{golang}
type Reading struct {
    DeviceID string `json:"device_id"`
}
func handle(w http.ResponseWriter,
    r *http.Request) {
    var in Reading
    json.NewDecoder(r.Body).Decode(&in)
    w.WriteHeader(http.StatusCreated)
    json.NewEncoder(w).Encode(in)
}
\end{golang}
    \end{column}
    \begin{column}{0.44\textwidth}
      {\bfseries Dart: the client side of the call}\par
\begin{dart}
final dio = Dio(BaseOptions(
  baseUrl: 'https://api.example.com',
));
final res = await dio.post(
  '/readings',
  data: {'device_id': id},
);
\end{dart}
    \end{column}
  \end{columns}
  \muted{\code{Reading.fromJson(res.data)} decodes the reply. The JSON tags are the contract between Go and Dart.}
  \note{Live demo: go run the file, curl the endpoint with a JSON body, then point the Flutter app at it. Total moving parts: one file and one binary. This is the whole backend a semester project usually needs.}
\end{frame}

\begin{frame}[eyebrow=Go backends]{Thundering herds and hedged requests}
  \begin{columns}[T]
    \begin{column}{0.62\textwidth}
      \begin{itemize}
        \item When a service fails briefly, ten thousand clients fail at once and every one retries after exactly one second. The retry wave is larger than the original load: a brief outage becomes a long one. That is a thundering herd.
        \item Full jitter spreads retries across the whole interval, so a recovering server sees load ramp up instead of arriving all at once.
        \item Hedging is the opposite problem: a small fraction of requests are simply slow. Send a second copy after a short delay, take whichever answers first, and cancel the other.
      \end{itemize}
    \end{column}
    \begin{column}{0.34\textwidth}
      {\usebeamerfont{small}\begin{tabular}{@{}p{16mm}p{25mm}@{}}
        \toprule
        \thead{Attempt} & \thead{Result} \\
        \midrule
        sent at 0 ms & silent at 250 ms, canceled \\
        sent at 250 ms & answers at 310 ms, used \\
        \bottomrule
      \end{tabular}}
      \vskip 2mm
      \muted{Cap it: hedge only past the 95th percentile, and only for safe or idempotent calls.}
    \end{column}
  \end{columns}
  \begin{takeaway}{Retries fix failures. Hedges fix slowness.}
    Both multiply load, so both need a cap.
  \end{takeaway}
  \note{The winner is whichever attempt replies first; the loser is simply canceled, not an error. Ask which methods are safe to hedge: GET freely, POST only behind an idempotency key the server remembers.}
\end{frame}

\divider{Part 4 \textperiodcentered\ API design for a mobile client}{Design decisions clients actually feel}[Pagination, versioning, conditional requests, compression, and rate limits]

\begin{frame}[fragile,eyebrow=API design]{Pagination with cursors, not page numbers}
  \begin{columns}[T]
    \begin{column}{0.56\textwidth}
      \begin{itemize}
        \item An offset-and-limit page breaks when rows are inserted or deleted between requests: page two can skip or repeat rows.
        \item A cursor is an opaque token that encodes the last item's sort key. The client sends it back to get the next page; nothing else has to stay stable.
        \item Cursors compose with real-time inserts, work with any sort order, and never need a total row count the database would have to compute.
      \end{itemize}
    \end{column}
    \begin{column}{0.4\textwidth}
\begin{golang}
// cursor, limit: query params
type Page struct {
    Items []Reading `json:"items"`
    Next  string `json:"next,omitempty"`
}
\end{golang}
    \end{column}
  \end{columns}
  \begin{takeaway}{A cursor is a fact about the last row returned.}
    It is not a page count the server has to keep consistent across writes.
  \end{takeaway}
  \note{Ask what happens to page two of an offset query if a row is deleted from page one while the user is still scrolling. Every item after it shifts up by one, and the user sees a duplicate or a skip.}
\end{frame}

\begin{frame}[eyebrow=API design]{Versioning: give clients time to update}
  \vskip 1mm
  \begin{itemize}
    \item A mobile client cannot be forced to update. Old app versions can stay installed for years, so the server must keep serving them.
    \item Put the version in the URL path, like \code{/v1/readings}, or in an Accept header. A path is simpler to log, cache, and route.
    \item Add fields freely: a client that ignores unknown JSON fields never breaks. Never remove or rename a field, and never change what a field means, without a new version.
    \item Deprecate loudly: a \code{Sunset} header and a fixed shutdown date, announced early and enforced only after clients had time to move.
  \end{itemize}
  \begin{takeaway}{Additive changes are free.}
    Anything else needs a new version number.
  \end{takeaway}
  \note{The failure mode students will hit in the project is renaming a JSON field mid-semester and breaking every teammate's app at once. A new version number, or an additive field, avoids that entirely.}
\end{frame}

\begin{frame}[eyebrow=API design]{ETags and conditional requests}
  \begin{itemize}
    \item An entity tag (ETag) is a short identifier for one exact representation of a resource: a hash or a version number.
    \item The server returns \code{ETag}. The client sends it back as \code{If-None-Match}; unchanged means 304 Not Modified, no body.
    \item \code{If-Match} on a write rejects it with 412 Precondition Failed if the resource changed since it was last read.
  \end{itemize}
  {\fontsize{8}{9.5}\selectfont\renewcommand\arraystretch{1.0}
  \begin{tabular}{@{}lll@{}}
    \toprule
    \thead{Header} & \thead{Direction} & \thead{Meaning} \\
    \midrule
    ETag & response & fingerprint of this representation \\
    If-None-Match & request & send a new copy only if this changed \\
    If-Match & request & write only if it still matches this version \\
    \bottomrule
  \end{tabular}}
  \begin{takeaway}{Conditional requests turn a download into a status code.}
    And a blind write into a safe one.
  \end{takeaway}
  \note{This is the same idea Lecture 9 used for optimistic concurrency in the sync engine, applied at the transport layer instead of inside the database.}
\end{frame}

\begin{frame}[eyebrow=API design]{Compression: fewer bytes on a metered radio}
  \vskip 1mm
  \begin{itemize}
    \item A phone often pays for every byte in time, battery, and sometimes money. Compression trades a little CPU time for fewer bytes on the radio.
    \item The client sends \code{Accept-Encoding: gzip}. The server compresses the body and answers with \code{Content-Encoding: gzip}. Most HTTP client libraries decompress transparently.
    \item JSON compresses well: repeated keys and whitespace shrink hard, often 70–90 percent smaller. A response already under about 1 KB rarely benefits enough to bother.
    \item Do not compress what is already compressed: images, video, and most binary formats gain nothing and waste CPU cycles.
  \end{itemize}
  \begin{takeaway}{Compress text responses by default.}
    Skip it for anything already binary.
  \end{takeaway}
  \note{Brotli, requested as br in Accept-Encoding, compresses text slightly better than gzip and is supported by most modern HTTP stacks, but gzip alone already covers most of the gain for a small API.}
\end{frame}

\begin{frame}[eyebrow=API design]{Rate limits: protecting the server from its clients}
  \vskip 1mm
  \begin{itemize}
    \item A rate limit caps how many requests one client, one API key, or one IP address may make in a window, so one runaway client or a retry storm cannot take the service down for everyone else.
    \item The token bucket is the common shape: a bucket refills at a fixed rate, and each request costs one token; a burst is allowed until the bucket is empty.
    \item The server answers 429 Too Many Requests with a \code{Retry-After} header, so a well-behaved client waits exactly that long instead of guessing.
    \item Return the budget in headers on every response, such as \code{X-RateLimit-Remaining}, so a client can back off before it gets rejected.
  \end{itemize}
  \begin{takeaway}{A limit the client can see in headers is one clients cooperate with.}
    Not one they fight.
  \end{takeaway}
  \note{Connect a rate limit back to the retry and hedge slides: a client that retries without jitter against a rate-limited endpoint produces the exact thundering herd from Part 3.}
\end{frame}

\divider{Part 5 \textperiodcentered\ Deploying and operating}{Shipping the binary and keeping it healthy}[A Docker image, Cloud Run, systemd, secrets, health checks, and logs]

\begin{frame}[fragile,eyebrow=Deploying]{A minimal Dockerfile for a Go binary}
\begin{codeblock}
FROM golang:1 AS build
WORKDIR /src
COPY . .
RUN CGO_ENABLED=0 go build -o /server ./cmd/server

FROM gcr.io/distroless/static
COPY --from=build /server /server
EXPOSE 8080
ENTRYPOINT ["/server"]
\end{codeblock}
  \vskip 2mm
  \muted{Multi-stage build: the first stage compiles with the full Go toolchain; the second stage copies only the binary into a distroless image with no shell and no package manager, which shrinks the image and the attack surface.}
  \note{Ask why the final image has no shell. It means an attacker who compromises the process cannot exec into a shell to look around, because there is no shell binary in the image at all.}
\end{frame}

\begin{frame}[fragile,eyebrow=Deploying]{The same image, deployed to Cloud Run}
  \begin{itemize}
    \item \code{gcloud run deploy} builds and deploys the container behind a managed load balancer. Nobody patches the machine.
    \item Concurrency, instances, memory, and CPU are flags, not hand-provisioned infrastructure.
    \item A new revision gets traffic gradually, so a bad deploy rolls back in seconds.
  \end{itemize}
\begin{shell}
gcloud run deploy api \
  --source . \
  --region europe-west1 \
  --set-env-vars ENV=production \
  --min-instances=0 --max-instances=10
\end{shell}
  \begin{takeaway}{Cloud Run is the serverless model from Part 2.}
    With your own container, not a vendor runtime.
  \end{takeaway}
  \note{This is the bridge between Part 2 and Part 3: the platform behavior (scale to zero, cold starts, per-request billing) is identical, only the runtime inside the box changed from a vendor function to a container you built.}
\end{frame}

\begin{frame}[fragile,eyebrow=Deploying]{A systemd unit on a VPS}
  \begin{columns}[T]
    \begin{column}{0.56\textwidth}
\begin{codeblock}
[Unit]
Description=api backend
After=network.target

[Service]
ExecStart=/usr/local/bin/server
Restart=always
RestartSec=2
User=api
EnvironmentFile=/etc/api/api.env

[Install]
WantedBy=multi-user.target
\end{codeblock}
    \end{column}
    \begin{column}{0.4\textwidth}
      \lead{One file, one service.}{\code{systemctl enable --now api} starts it and keeps it running across reboots.}
      \muted{\code{Restart=always} with a short \code{RestartSec} means a crash recovers in about two seconds, with nobody paged at 3 a.m.}
    \end{column}
  \end{columns}
  \note{Point out that this is the whole operational story for a project backend: one unit file, one restart policy, one environment file. Nothing here needs Kubernetes.}
\end{frame}

\begin{frame}[fragile,eyebrow=Deploying]{Environment variables and secrets}
  \begin{columns}[T]
    \begin{column}{0.58\textwidth}
      \begin{itemize}
        \item Configuration that changes between environments belongs in environment variables, read once at startup, never hardcoded.
        \item Secrets are configuration too, but never in the code or the git repository. A leaked key is compromised the moment it is pushed, even after deletion.
        \item On a VPS: a root-only environment file loaded by systemd. On Cloud Run: a secret manager, never baked into the image.
      \end{itemize}
    \end{column}
    \begin{column}{0.38\textwidth}
\begin{golang}
dsn := os.Getenv("DATABASE_URL")
if dsn == "" {
    log.Fatal("DATABASE_URL is required")
}
\end{golang}
    \end{column}
  \end{columns}
  \begin{takeaway}{A required variable that is missing should crash at startup.}
    Not on the first request, hours later.
  \end{takeaway}
  \note{Ask the room whether they have ever committed an API key by accident. Most have. The fix is not remembering harder: it is a .gitignore entry and a pre-commit check that scans for secret-shaped strings.}
\end{frame}

\begin{frame}[fragile,eyebrow=Deploying]{Health checks: liveness and readiness}
  \vskip 1mm
  \begin{itemize}
    \item Liveness answers one question: is the process still running and not deadlocked. A failing liveness check gets the process restarted, so keep it cheap and dependency-free.
    \item Readiness answers a different question: is this instance ready for traffic. Failing it removes the instance from the load balancer without a restart.
    \item Cloud Run, most orchestrators, and load balancers expect a plain HTTP endpoint that returns 200 when healthy, checked every few seconds.
  \end{itemize}
\begin{golang}
http.HandleFunc("/healthz", func(w http.ResponseWriter, r *http.Request) {
    w.WriteHeader(http.StatusOK)
})
\end{golang}
  \begin{takeaway}{Liveness restarts the process. Readiness redirects traffic.}
    Confusing the two either restarts a healthy process or sends traffic to a broken one.
  \end{takeaway}
  \note{The distinction matters most under database failure: the process itself is fine, so liveness should stay green while readiness goes red and the load balancer quietly stops sending it traffic.}
\end{frame}

\begin{frame}[fragile,eyebrow=Deploying]{Structured logs: fields, not sentences}
  \begin{columns}[T]
    \begin{column}{0.58\textwidth}
      \begin{itemize}
        \item A log line meant for a person cannot be filtered, aggregated, or alerted on reliably. A structured log line is a set of key-value fields a machine can query.
        \item Go's log/slog, in the standard library, writes JSON through a JSON handler: a level, a message, and any fields attached, one line per event.
        \item Always include a request identifier a client sends or the server generates, so one failing mobile session can be traced through every log line it touched.
      \end{itemize}
    \end{column}
    \begin{column}{0.38\textwidth}
\begin{golang}
h := slog.NewJSONHandler(
  os.Stdout, nil,
)
logger := slog.New(h)
logger.Info("reading saved",
  "device_id", id,
  "duration_ms", ms)
\end{golang}
    \end{column}
  \end{columns}
  \begin{takeaway}{If you cannot turn it into a chart, it is not a structured log yet.}
    A sentence in a log file is for a human reading it live, once.
  \end{takeaway}
  \note{Log the boring successful requests too, not just the errors: status code, duration, and device identifier on every call is what tells you the baseline changed before it becomes an outage.}
\end{frame}

\begin{frame}[eyebrow=Deploying]{What the project backend needs}
  \vskip 1mm
  \begin{itemize}
    \item Lab 5's sync server is this lecture in miniature: a small Go program, an in-memory list with a version per row, one endpoint that accepts offline writes and returns what changed since a cursor.
    \item Lab 6 adds gRPC (Lecture 11) to the same binary: one unary call and one server-streaming call, both served from the process that already answers HTTP.
    \item Keep it on one small VPS or one Cloud Run service for the whole semester. The project needs a reliable box, not a distributed system.
    \item Everything in this lecture applies at project scale: environment variables for broker and database URLs, a health check the grading script can poll, and structured logs for the night before the defense.
  \end{itemize}
  \begin{takeaway}{The project backend is small on purpose.}
    The skill being graded is running it correctly, not scaling it.
  \end{takeaway}
  \note{Reassure the room: nobody needs Kubernetes for this project. One VPS, one binary, one systemd unit, and the checklist from this lecture is a complete, defensible backend.}
\end{frame}

\begin{frame}[eyebrow=Recap]{Three things to remember}
  \vskip 2mm
  \begin{enumerate}
    \item The client is untrusted input. \muted{Every rule that matters is checked again on the server.}
    \item Serverless bills per request and cold-starts; a VPS bills per month and never does. \muted{Pick by cost shape, not by fashion.}
    \item A cursor, an ETag, and a rate-limit header are what make an API pleasant to build a mobile client against. \muted{Small, additive changes keep old app versions working.}
  \end{enumerate}
  \vskip 5mm
  \kv{Further reading}{\link{https://go.dev/doc/tutorial/web-service-gin}{go.dev/doc/tutorial/web-service-gin} \textperiodcentered\ \link{https://pkg.go.dev/net/http}{pkg.go.dev/net/http} \textperiodcentered\ \link{https://pkg.go.dev/log/slog}{pkg.go.dev/log/slog} \textperiodcentered\ \link{https://cloud.google.com/run/docs}{cloud.google.com/run/docs}}
  \note{Close by pointing at the Lab 5 branch due next: the sync server they build there is a direct application of the Go backend and deployment material from this lecture.}
\end{frame}

\closingframe{Questions?}{dinu\_stefan.rusu@upb.ro \textperiodcentered\ Microsoft Teams: course channel}

\end{document}
